\frametitle{Two Capacity Diagnostics, Reported Beside Every Number}
\footnotesize

A DeepONet output has the bilinear form $\delta e=\sum_k b_k(r)T_k(x)$, so the correction
lies in $\operatorname{span}(T)$ \emph{whatever the branch returns}. A small $\muF$ can
therefore be a statement about the trunk rather than about smoothing --- and the first run
of this experiment measured exactly that. \alert{The floor is computable before training.}

\begin{columns}[T]
\begin{column}{0.52\textwidth}
\centering
\textbf{1. The trunk sets a floor} \quad
$\operatorname{reach}_F(e)=\bigl\|\operatorname{proj}_{\operatorname{span}(T)}e_F
\bigr\|_A^2/\normA{e_F}^2$, \ $\muF\ge\mathbb{E}\bigl[\sqrt{1-\operatorname{reach}_F}\bigr]$.

\vspace{1mm}
{\scriptsize
\begin{tabular}{@{}lrrr@{}}
\toprule
Trunk & dim & rank & floor on $\muF$ \\
\midrule
Fourier $K=2$ & 48 & 24 & $0.997784$ \\
Fourier $K=6$ & 336 & 168 & $0.965419$ \\
Fourier $K=10$ & 880 & 440 & $0.786705$ \\
Fourier $K=12$ & 1248 & 624 & $0.563591$ \\
Fourier $K=14$ & 1680 & 840 & $0.286669$ \\
\textbf{Fourier $K=16$} & 2176 & 1023 & \textbf{$0.000675$} \\
\bottomrule
\end{tabular}}

\vspace{1mm}
\raggedright
The reason is structural, not numerical. $\operatorname{range}(P)$ is what the coarse
lattice can represent ($16\times16$, Nyquist $8$), so its $A$-orthogonal complement is
essentially the band \emph{between the coarse and the fine Nyquist} --- modes
$8\lesssim|k|\le16$. A Fourier trunk needs modes up to $16$. The headline run therefore
\emph{freezes} the trunk at $K=16$: then $\operatorname{span}(T)=\R^n$, $T^{\mathsf T}T=I$,
and the floor is \alert{zero} --- provably non-binding.
\end{column}
\begin{column}{0.45\textwidth}
\centering
\textbf{2. The branch sets a second floor}

\raggedright
The branch narrows to \emph{width} units, so its Jacobian factors through
$\R^{\text{width}}$, while the target of this stage is a \emph{linear map of rank $768$}.
A branch narrower than that cannot express the exact correction, whatever the optimiser
does. The reference curve is the best \emph{linear} map of rank $\le k$ fitted to the same
pairs with the same weighting --- a property of the \emph{data}, no network is trained.

\vspace{1mm}
{\scriptsize
\begin{tabular}{@{}rrrr@{}}
\toprule
rank $k$ & $\muF$ & $\muC$ & $\muC/\muF$ \\
\midrule
64 & $0.8057$ & $1.0103$ & $1.25$ \\
128 & $0.7160$ & $1.0150$ & $1.42$ \\
\textbf{256} & \textbf{$0.5860$} & $1.0200$ & $1.74$ \\
384 & $0.4768$ & $1.0224$ & $2.14$ \\
512 & $0.3666$ & $1.0237$ & $2.79$ \\
640 & $0.2375$ & $1.0245$ & $4.31$ \\
\textbf{768} & \textbf{$0.0000$} & $1.0249$ & $\infty$ \\
\bottomrule
\end{tabular}}

\vspace{1mm}
\raggedright
The curve saturates at the complement dimension --- the rank a branch \emph{must} have to
do the job. And $\muC\approx1.02$ at \emph{every} rank: selectivity is not something the
network has to discover. The interesting question is therefore not whether selectivity
appears but how well the complement is reduced --- and a branch of width $w$ sits near the
$k=w$ row.
\end{column}
\end{columns}
